\documentclass{ltjsarticle}

\setlength\parindent{0pt}
\setlength{\columnseprule}{0.4pt}

\newcommand{\sk}{\vskip\baselineskip}

\usepackage[margin=15mm]{geometry}
\usepackage{amsmath}
\usepackage{amssymb}
\usepackage{amsfonts}
\usepackage{mathrsfs}
\usepackage{bm}
\usepackage{here}
\usepackage{bbm}
\usepackage{graphicx}
\usepackage[unicode]{hyperref}
\usepackage{mathtools}
\usepackage{url}
\usepackage{ascmac}
\usepackage{physics}
\usepackage[version=4]{mhchem}
\usepackage{listings}

\lstset{
language=Python, 
basicstyle={\small},
identifierstyle={\small},
commentstyle={\small\itshape},
keywordstyle={\small\bfseries},
ndkeywordstyle={\small},
stringstyle={\small\ttfamily},
frame={tblr},
breaklines=true,
columns=[l]{fullflexible},
xrightmargin=0zw,
xleftmargin=3zw,
numberstyle={\scriptsize},
stepnumber=1,
numbersep=1zw,
lineskip=-0.5ex
}

\mathtoolsset{showonlyrefs=true}

\title{ランダム行列 常識集}
\author{東京大学 理学系研究科物理学専攻 修士1年 平松大門}
\begin{document}
\maketitle
\section{Getting Started}
\begin{itemize}
    \item PDF(Probabilty Density Function)とは？
    
    確率密度関数$p(x)$のこと

    \sk

    \item JPDF(Joint Probability Density Function)とは？
    
    同時確率密度関数$p(x_1,\cdots x_n)$のこと

    \sk
    \item CDF(Cumulative Distribution Function)とは？
    
    累積分布関数$F(x):=\int_{-\infty}^{x}dt \,p(t)$のこと
    
    \sk
    \item Gaussian Ensemblesの代表例3つは？
    \begin{table}[h]
\centering
\begin{tabular}{|c|c|c|c|}
\hline
略称 & 正式名称&行列の性質&Dyson Index $\beta$\\
\hline
GOE & Gaussian Orthogonal Ensemble & 実対称行列\quad $H=H^T$ &1
\\
\hline
GUE & Gaussian Unitary Ensemble&エルミート行列 \quad$H=H^\dagger$ &2
\\
\hline
GSE &Gaussian Symplectic Ensemble& 四元数自己双対エルミート行列 &4
\\
\hline
\end{tabular}
\end{table}

名前は分布が不変となる変換に由来する。
Dyson Indexは行列の成分が持つ実自由度の数を表す。

\sk
\item 重積分の変数変換は？
\[
\int\prod_i dx_i \,f(x_1, \cdots ,x_n) = \int \prod_i dy_i\, \left|\frac{\partial (x_1, \cdots , x_n)}{\partial (y_1, \cdots ,y_n)}\right|f\left(x_1(\bm{y}), \cdots, x_n(\bm{y})\right)
\]
$d\bm{y}$だけ動いた時に$\bm{x}$空間で動く体積が$\left|\frac{\partial (x_1, \cdots , x_n)}{\partial (y_1, \cdots ,y_n)}\right|$

\end{itemize}

\newpage

\section{Value the Eigenvalue}
\subsection{Appetizer: Wigner's surmise}
\begin{itemize}
    \item $2\times 2$のGOE行列$H_s := \begin{pmatrix}
        x_1 & x_3\\
        x_3 & x_2
    \end{pmatrix}$
    で、$x_1, x_2\sim N(0, 1)$、$x_3\sim N(0,\frac{1}{2})$に従う時、
    固有値間隔$s:=|\lambda_1 -\lambda_2|$の分布$p(s)$はどのようになるか？

    \[
    p(s) = \frac{s}{2}e^{-\frac{s^2}{4}}
    \]
    この分布は固有値同士が重なる確率が0、つまり固有値同士は反発することを意味する。

    証明は、$p(s)=\int dx_1 \frac{e^{-\frac{x_1^2}{2}}}{\sqrt{2\pi}}\,dx_2  \frac{e^{-\frac{x_2^2}{2}}}{\sqrt{2\pi}}\,dx_3\frac{e^{-x_3^2}}{\sqrt{\pi}}\,\delta(s-|\lambda_1 - \lambda_2|)$を変数変換して計算するだけ。
    変数変換したあと積分が実行できるためには、非対角成分の分散が対角成分の分散の$\frac{1}{2}$になっていることが必要。
    
    \sk
    \item 確率分布$p(x)$を期待値が1になるようリスケールした分布$\tilde{p}(x)$は？
    
    \[
    \tilde{p}(x) := \langle x \rangle \,p(\langle x \rangle x)
    \]
    横向きに$\frac{1}{\langle x\rangle}$倍、縦向きに$\langle x \rangle$倍すれば良い。

    \sk
    \item Wigner予想とは？
    
    ランダム行列の固有値間隔の分布についての予想。$p(s)$を$s$の期待値が1になるようリスケールした分布$\tilde{p}(s)$は
    \[
    \tilde{p}(s)=\frac{\pi}{2}s\,\exp(-\frac{\pi}{4}s^2)
    \]
    となるが、$N\times N$でもリスケールしたものは似た分布になると予想した。
\end{itemize}

\subsection{Eigenvalues as correlated random variables}
\begin{itemize}
    \item level repulsionとは？
    
    ランダム行列の固有値同士は反発しあい、独立ではないこと
\end{itemize}

\sk
\subsection{Compare with the Spacings Between i.i.d.'s}
\begin{itemize}
    \item i.i.d.(independent and indentically disributed)とは?
    
    確率変数がそれぞれ同じ確率分布に従って、かつ互いに独立している状況のこと。

    \sk
    \item i.i.d.である確率変数の間隔分布は？\\
    地点$x$の近傍においてある変数とその次の変数までの距離が$s$となる確率密度$p(s)$は
    \begin{align}
        p(s)=Np(x)\,e^{-Np(x)s}
    \end{align}
    i.i.d.では近傍に含まれる変数の個数が2項分布に従い、$N\to\infty$においてそれがPoisson分布に変化するため
    
    \sk
    \item $p(x)dx$を含む積分をする時は？\\
    $p(x)$のCDFである$F(x)$が現れたとしても$dF(x)=p(x)dx$となるので、積分対象にCDFが現れても良い
\end{itemize}

\sk
\subsection{Jpdf of Eigenvalues of Gaussian Matrices}
\begin{itemize}
    \item $N(0,1)$に従うGaussian Ensembleの固有値$\{\lambda_i\}$のJPDFは？\\
    \begin{align}
        \rho(\lambda_1,\cdots,\lambda_N)=\frac{1}{Z_{N,\beta}}\exp\left(-\frac{1}{2}\displaystyle\sum_i \lambda_i^2\right) \prod_{j<k}\|\lambda_j-\lambda_k\|^\beta
    \end{align}
ただし、$Z_{N,\beta}=(2\pi)^{\frac{N}{2}}\displaystyle\prod_j\frac{\Gamma(1+\frac{j\beta}{2})}{\Gamma(1+\frac{\beta}{2})}$
\end{itemize}

\newpage
\section{Classified Material}

\newpage
\section{The Fluid Semicircle}

\newpage
\section{Suddle-Point-of-View}


\newpage
\section{Time for a Change}

\newpage
\section{Meet Vandermonde}

\newpage
\section{Resolve(nt) Semicircle}

\newpage
\section{One Pager on Eigenvectors}
\begin{itemize}
    \item $N\times N$のGUE Ensembleの規格化された固有ベクトルの1成分$c\in \mathbb{C}$の2乗ノルム$y:=\|c\|^2$の分布$P_{\mathrm{GUE}}(y)$は？
    \begin{align}
        P_{\mathrm{GUE}}(y)=(N-1)(1-y)^{N-2}
    \end{align}
    固有ベクトルの成分は0付近に集中している。
    \sk
    \item GUEでの$y$の期待値は？
    \begin{align}
        E[y] = \frac{1}{N}
    \end{align}

    \sk
     \item GUEでの$y$の分散は？
     \begin{align}
        V[y] = \frac{1}{N^2}\left(1-\frac{1}{N+1}\right)
     \end{align}

    \sk
    \item GOEでの$y$の期待値は？
\end{itemize}

\newpage
\section{Finite N}

\end{document}